\documentclass[10pt,a4paper]{amsart}
\usepackage[margin=19mm]{geometry}
\usepackage{amsmath,amssymb,mathtools}
\usepackage[scr=rsfs,cal=euler]{mathalfa}
\usepackage{xfrac}
\usepackage[unicode,hidelinks,bookmarks=false]{hyperref}
\newcommand{\ie}{i.e.\ }
\newcommand{\cf}{cf.\ }
\newcommand{\resp}{respectively\ }
\newcommand{\Subsection}{\subsection}
\providecommand{\nobreakdash}{\nobreak\hbox{-}}
\setlength{\emergencystretch}{2em}
\multlinegap0pt
\usepackage{bm}
\usepackage{bbm}
\usepackage{dsfont}
\usepackage{xspace}
\usepackage{cancel}
\usepackage{mathtools}
\usepackage{amsthm}
\makeatletter
\@ifundefined{definition}{\newtheorem{definition}{Definition}}{}
\@ifundefined{proposition}{\newtheorem{proposition}{Proposition}}{}
\makeatother
\title[Segmentation-Based Regression for Quantum Neural Networks]{Segmentation-Based Regression for Quantum Neural Networks}
\author{James C. Hateley}
\date{Source: arXiv:2507.00065v1 (CC BY 4.0)}
\begin{document}
\maketitle

\noindent\emph{Source and licence.} Segmentation-Based Regression for Quantum Neural Networks. Version arXiv:2507.00065v1 (CC BY 4.0), distributed under the Creative Commons Attribution 4.0 International licence, as recorded in the arXiv licence metadata for this exact version and shown on the abstract page https://arxiv.org/abs/2507.00065v1. Full rights notice: licenses/arxiv-2507.00065-CC-BY-4.0.txt.\par\smallskip

\noindent\textit{Editorial scope.} Counted unit: the Proposition whose source label is \texttt{prop:clipping} in the article (source0.tex lines 671--687, source label \texttt{prop:clipping}), delivered together with the article's own complete original proof of it (source0.tex lines 689--717, one \texttt{proof} environment of 29 lines). No mathematics is written, repaired or re-derived: statement and proof are the article's own text line for line. Required context retained uncounted: def:digit\_representation (lines 77--82); def:signed\_digits (lines 105--111). Every definition, display and label of those lines is retained unchanged. Omitted from this package and not redistributed: 1--76, 83--104, 112--670, 688--688, 718--1266. Nothing inside a retained excerpt is omitted or altered: every retained body is delivered exactly as the article's lines are written, so no omission receipt of any kind accompanies this package. Citations: the retained excerpts carry no citation command at all, so no bibliography entry of the article is transcribed and no reference is redistributed. Reference adaptation: none was needed and none was made; every reference inside the delivered unit resolves inside it, so no unresolved reference or citation is produced. Printed slips kept literal: no printed word, symbol, hypothesis or number of the counted statement or of its proof was altered. Layout and class: the article's own class is \texttt{revtex4} with packages including amsmath, amsthm, physics, mathrsfs, amsfonts, amssymb, graphicx, float, babel, hyperref, enumerate; the wrapper is the lane's pinned \texttt{amsart} editorial wrapper at 10pt on A4, which restates the theorem-like environments the retained text needs and adds the article's own macro definitions that the retained text needs (none), with extra wrapper packages (bm, bbm, dsfont, xspace, cancel, mathtools). The delivered numbering is the unit's own. Assets: no retained span contains a figure environment, a graphics-inclusion command, a PDF-page inclusion, a table or a TikZ picture; no image file is copied into this package, the package declares no asset dependency and no image file is redistributed with it. Licence: version arXiv:2507.00065v1 of this article is distributed under the Creative Commons Attribution 4.0 International licence, as recorded in the arXiv licence metadata for this exact version and shown on the abstract page https://arxiv.org/abs/2507.00065v1. A full rights notice is delivered at licenses/arxiv-2507.00065-CC-BY-4.0.txt. Characters: every retained excerpt is pure ASCII, so the delivered engine, XeLaTeX with its default Latin Modern text font, reports no missing character.
\begin{definition}[Base-$b$ Digit Expansion]\label{def:digit_representation}
Let $b \in \mathbb{Z}_{\geq 2}$ be a fixed digit base, and $(n,m) \in \mathbb{N}^2$ the digit resolution, with $n$ digits left of the radix point and $m$ to the right. A real number $y \in \mathbb{R}$ admits a base-$b$ expansion if
\begin{equation}
y = \sum_{i=-m}^{n} y_i b^i, \quad y_i \in \{0,1,\ldots,b-1\}.
\end{equation}
\end{definition}

\begin{definition}[Signed Digit Expansion]\label{def:signed_digits}
In a signed expansion, each digit $y_i$ belongs to the symmetric set $\{-\lfloor b/2 \rfloor, \ldots, \lfloor (b-1)/2 \rfloor\}$. The representable range becomes symmetric about zero:
\begin{equation}
\theta_{\min}^{\text{signed}} = -\sum_{i=-m}^{n} \left\lfloor \frac{b}{2} \right\rfloor b^i, \quad
\theta_{\max}^{\text{signed}} = \sum_{i=-m}^{n} \left\lfloor \frac{b-1}{2} \right\rfloor b^i.
\end{equation}
\end{definition}

\begin{proposition}[Clipping Error Due to Digit Range Constraints]
\label{prop:clipping}
Let \( \theta^* \in \mathbb{R} \) be a true target value, and suppose the digitized reconstruction \( \tilde{\theta} \in Y_{b,n,m} \) is constrained to lie within the representable range determined by an \emph{unsigned} base-\( b \) digit expansion (Definition~\ref{def:digit_representation}). In this case, the minimum and maximum representable values are given by: $\theta_{\min} := 0$,
\[
\theta_{\max} := \sum_{i=-m}^{n} (b-1) b^i = (b^{n+1} - b^{-m})(1 - b^{-1}).
\]
The following bound holds on the worst-case clipping error due to range saturation:
\begin{equation}
|\tilde{\theta} - \theta^*| \leq \max \big\{ 0,\, |\theta^* - \theta_{\max}|,\, |\theta^* - \theta_{\min}| \big\}.
\end{equation}
In particular, if \( \theta^* > \theta_{\max} \), then the recovered value saturates at \( \tilde{\theta} = \theta_{\max} \), and the reconstruction error satisfies \( |\tilde{\theta} - \theta^*| = |\theta^* - \theta_{\max}| \).

\medskip

\noindent
\textit{Remark.} For signed digit expansions (Definition~\ref{def:signed_digits}), the representable interval becomes symmetric and centered around zero. An analogous bound applies, with adjusted endpoints \( \theta_{\min}^{\mathrm{(signed)}}, \theta_{\max}^{\mathrm{(signed)}} \) defined accordingly.
\end{proposition}

\begin{proof}
Let \( \theta^* \in \mathbb{R} \) denote the true target value, and let \( \tilde{\theta} \in Y_{b,n,m} \) be the closest representable value to \( \theta^* \) within the digitized lattice. Since the lattice is defined by fixed digit resolution and non-negative digit values, all representable quantities lie within the interval
\begin{equation}
[\theta_{\min}, \theta_{\max}] = \left[ 0, \sum_{i=-m}^{n} (b-1) b^i \right].
\end{equation}
By definition, if \( \theta^* \in [\theta_{\min}, \theta_{\max}] \), then there exists a digitized approximation \( \tilde{\theta} \in Y_{b,n,m} \) such that \( |\tilde{\theta} - \theta^*| \leq \delta \), where \( \delta \) is the quantization resolution (e.g., \( \delta \sim b^{-m} \)). In this case, no clipping occurs and the error is bounded by the quantization error.

However, if \( \theta^* \notin [\theta_{\min}, \theta_{\max}] \), then the closest representable value is given by the boundary point of the interval:
\[
\tilde{\theta} = 
\begin{cases}
\theta_{\min} & \text{if } \theta^* < \theta_{\min}, \\
\theta_{\max} & \text{if } \theta^* > \theta_{\max}.
\end{cases}
\]
Hence, the reconstruction error is
\[
|\tilde{\theta} - \theta^*| = 
\begin{cases}
|\theta^* - \theta_{\min}| & \text{if } \theta^* < \theta_{\min}, \\
|\theta^* - \theta_{\max}| & \text{if } \theta^* > \theta_{\max}.
\end{cases}
\]
Combining all cases, we obtain the uniform bound
\begin{equation}
|\tilde{\theta} - \theta^*| \leq \max \left\{ 0,\, |\theta^* - \theta_{\max}|,\, |\theta^* - \theta_{\min}| \right\}.
\end{equation}
This concludes the proof.
\end{proof}

\end{document}
